\documentclass[a4paper]{article}

\usepackage{graphicx}
\usepackage{amsmath,amssymb}
\usepackage{minted}
\usepackage{multirow}
\usepackage{float}
\usepackage{algorithm}
\usepackage{algpseudocode}
\usepackage{todonotes}
\usepackage{forest}
\usepackage{xstring}
\usepackage{xcolor}
\usepackage[most]{tcolorbox}
\usepackage{xparse}
\usepackage{natbib}
\usepackage{adjustbox}

\usetikzlibrary{graphs,graphdrawing}
\usegdlibrary{layered}

% for external graphviz
\input{/home/donkarlo/Dropbox/repo/nd_language_ai_project/src/nd_language_ai/formal/computer/programming/tex/latex/code_clips/external_graphviz.tex}

% for tags
\input{/home/donkarlo/Dropbox/repo/nd_language_ai_project/src/nd_language_ai/formal/computer/programming/tex/latex/part/control_sequence/kind/semtag/semtag.tex}

% for links
\input{/home/donkarlo/Dropbox/repo/nd_language_ai_project/src/nd_language_ai/formal/computer/programming/tex/latex/code_clips/seehere.tex}

% for nested lists and sections
\input{/home/donkarlo/Dropbox/repo/nd_language_ai_project/src/nd_language_ai/formal/computer/programming/tex/latex/code_clips/deep_structure.tex}

\title{Mahalanobis Distance}
\date{}
\author{Mohammad Rahmani}

\begin{document}
   \maketitle

\section{Scoring rule}
For a feature vector \(x\), reference mean \(\mu\), and invertible reference covariance matrix \(\Sigma\), the squared distance is
\[
D^2(x)=(x-\mu)^\top\Sigma^{-1}(x-\mu).
\]
Covariance accounts for feature scale and correlation. Departures along directions with low reference variance receive more weight.

\section{Example}
Suppose standardized revenue and order count have correlation \(0.9\). The point \((1,-1)\) has \(D^2=20\), while \((1,1)\) has \(D^2\approx 1.05\). Each point has the same Euclidean distance from the mean.

\section{Practical considerations}
Under a multivariate Gaussian model with known parameters, \(D^2\) follows a chi-square distribution with degrees of freedom equal to the feature count. With estimated parameters, check calibration on held-out reference data. Regularization helps with poorly conditioned covariance, and robust covariance estimation helps when the reference sample contains outliers.



   \bibliographystyle{plainnat}
   \bibliography{/home/donkarlo/Dropbox/repo/nd_shared_research/bibliography/bibliography.bib}
\end{document}
